\documentclass[a4paper,12pt]{article}
\usepackage[T2A]{fontenc}
\usepackage[utf8]{inputenc}
\usepackage[russian]{babel}
\usepackage{amsmath,amssymb,icomma}
\usepackage[a4paper,margin=20mm,headheight=25pt,headsep=8mm,footskip=12mm]{geometry}
\usepackage{microtype,booktabs,tabularx,enumitem,fancyhdr,titlesec,setspace,xcolor,tikz}
\usetikzlibrary{arrows.meta,positioning,shapes.geometric,calc}
\definecolor{accent}{HTML}{008080}
\definecolor{darkaccent}{HTML}{004D4D}
\definecolor{textgray}{HTML}{333333}
\usepackage[hidelinks]{hyperref}
\color{textgray}
\setlength{\parindent}{0pt}
\setlength{\parskip}{6pt}
\setstretch{1.22}
\setlist{nosep,leftmargin=*,itemsep=4pt}
\titleformat{\section}{\Large\normalfont\color{darkaccent}}{\thesection.}{0.5em}{}[\vspace{2mm}\color{accent}\titlerule]
\titleformat{\subsection}{\large\normalfont\color{darkaccent}}{}{0pt}{}
\titlespacing*{\section}{0pt}{0pt}{10pt}
\titlespacing*{\subsection}{0pt}{10pt}{4pt}
\pagestyle{fancy}\fancyhf{}
\fancyhead[L]{\small Динамика. Второй закон Ньютона}
\fancyhead[R]{\small Филипп Лузянин}
\fancyfoot[C]{\small\thepage}
\renewcommand{\headrulewidth}{0.4pt}
\makeatletter\renewcommand{\headrule}{\hbox to\headwidth{\color{accent}\leaders\hrule height \headrulewidth\hfill}}\makeatother
\newcommand{\unit}[1]{\,\text{#1}}
\newcommand{\fr}{F_{\text{тр}}}
\newcommand{\Ft}{F_{\text{т}}}
\newcommand{\example}[1]{\par\medskip\textcolor{darkaccent}{\textbf{#1}}\par}
\tikzset{force/.style={-{Stealth[length=2.5mm]},accent,thick},axis/.style={-{Stealth[length=2mm]},gray},body/.style={draw=textgray,fill=white,thick},motion/.style={-{Stealth[length=2.5mm]},darkaccent,dashed,thick},flow/.style={draw=accent,rounded corners=2pt,align=center,text width=10.8cm,minimum height=1.05cm,inner sep=5pt},terminal/.style={ellipse,draw=accent,align=center,inner sep=5pt},link/.style={-{Stealth[length=2.5mm]},gray,thick}}
\begin{document}
\begin{titlepage}
\centering\thispagestyle{empty}
\vspace*{28mm}
{\Large\color{darkaccent} ФИЗИКА\par}
\vspace{16mm}
{\huge Чек-лист\par}
\vspace{12mm}
{\LARGE\color{accent} Динамика: второй закон Ньютона,\\ силы трения и связанные тела\par}
\vspace{18mm}
{\large 9 класс. Школьная программа\par}
\vfill
{\large Филипп Лузянин\par}
\vspace{12mm}
25 ноября 2026 г.\par
\vspace{19mm}
{\small Лёвин Артём Александрович\\ эксклюзивно для Филиппа Лузянина\par}
\vspace{18mm}
\begin{tikzpicture}\draw[accent,line width=.65pt] (0,0) .. controls (2,0.85) and (4,-0.85) .. (6,0) .. controls (8,0.85) and (10,-0.85) .. (12,0);\end{tikzpicture}
\vspace{5mm}
\end{titlepage}
\setcounter{page}{1}
\section{Силы и второй закон Ньютона}
\textbf{Главная идея:} ускорение тела определяется равнодействующей всех сил, приложенных к этому телу. В инерциальной системе отсчёта
\[\sum\vec F=m\vec a,\qquad \sum F_x=ma_x,\quad \sum F_y=ma_y.\]
Здесь $m$ --- масса (кг), $a$ --- ускорение (м/с$^2$), сила измеряется в ньютонах: $1\unit{Н}=1\unit{кг}\cdot\text{м}/\text{с}^2$.
\begin{center}
\renewcommand{\arraystretch}{1.3}
\begin{tabularx}{\linewidth}{@{}l l X@{}}\toprule
Сила & Обозначение & Направление и особенности\\\midrule
Сила тяжести & $m\vec g$ & Вертикально вниз; $F_{\text{тяж}}=mg$.\\
Реакция опоры & $\vec N$ & Перпендикулярно поверхности, от опоры к телу.\\
Трение скольжения & $\vec F_{\text{тр}}$ & Против движения тела относительно поверхности.\\
Натяжение нити & $\vec T$ & Вдоль натянутой нити, от тела.\\
Сила тяги & $\vec F_{\text{т}}$ & В направлении действия тяги.\\\bottomrule
\end{tabularx}
\end{center}
\subsection{Сила трения скольжения}
\[\fr=\mu N.\]
$\mu$ --- безразмерный коэффициент трения скольжения, $\mu\geq0$. В обычных школьных задачах часто $\mu<1$, но универсального ограничения сверху единицей нет. Для трения покоя модуль определяется условиями равновесия и не превышает $\mu_{\text{п}}N$.
\subsection{Как выбирать знаки проекций}
Если сила направлена вдоль положительного направления оси, её проекция положительна; если против оси --- отрицательна; если перпендикулярно оси --- равна нулю. При наклоне силы к оси под углом $\alpha$ её проекция на эту ось равна $F\cos\alpha$ (угол отсчитывается от положительного направления оси).
\textbf{Важно:} скорость и ускорение могут быть направлены по-разному. Силу трения определяют по направлению скольжения. На схемах длины стрелок условны.
\clearpage
\section{Движение по горизонтальной поверхности}
Тело скользит вправо по горизонтальной опоре. Сила тяги горизонтальна, вертикального ускорения нет. Ось $Ox$ направим вправо, $Oy$ --- вверх.
\begin{center}
\begin{tikzpicture}[scale=1.05]
\draw[gray] (-3.8,-.52)--(3.8,-.52);
\draw[axis] (-3.5,0)--(3.7,0) node[right]{$x$};
\draw[axis] (0,-1.55)--(0,2) node[above]{$y$};
\draw[body] (-.65,-.5) rectangle (.65,.5);
\draw[force] (0,0)--(2.45,0) node[above]{$\vec F_{\text{т}}$};
\draw[force] (0,0)--(-1.65,0) node[above]{$\vec F_{\text{тр}}$};
\draw[force] (0,0)--(0,1.5) node[right]{$\vec N$};
\draw[force] (0,0)--(0,-1.35) node[right]{$m\vec g$};
\draw[motion] (1.15,1.35)--(2.85,1.35) node[right]{$\vec v$};
\end{tikzpicture}
\end{center}
Второй закон Ньютона в проекциях:
\[\begin{aligned} Ox:\quad &\Ft-\fr=ma_x,\\ Oy:\quad &N-mg=0.\end{aligned}\]
Из второго уравнения $N=mg$, поэтому $\fr=\mu mg$. Подстановка в первое даёт
\[\boxed{a_x=\frac{\Ft-\mu mg}{m}=\frac{\Ft}{m}-\mu g.}\]
При $\Ft>\mu mg$ скорость скользящего вправо тела увеличивается; при равенстве --- постоянна; при $\Ft<\mu mg$ тело замедляется до остановки.
\example{Пример. Найдите ускорение.}
Масса бруска $5\unit{кг}$, горизонтальная сила тяги $30\unit{Н}$, $\mu=0,2$, $g=10\unit{м/с}^2$.
\[N=50\unit{Н},\quad \fr=0,2\cdot50=10\unit{Н},\]
\[a_x=\frac{30-10}{5}=4\unit{м/с}^2.\]
Ускорение направлено вправо. Если бы тяга составляла $5\unit{Н}$ при том же скольжении вправо, то $a_x=-1\unit{м/с}^2$, то есть ускорение было бы направлено влево.
\clearpage
\section{Два тела, связанные нитью}
Два бруска скользят вправо по горизонтальному столу. К правому бруску приложена горизонтальная сила $F$. Нить лёгкая, нерастяжимая, натянутая; ускорения брусков одинаковы, модули сил натяжения на её концах равны. Коэффициенты трения могут различаться; для гладкой опоры соответствующий коэффициент равен нулю.
\begin{center}
\begin{tikzpicture}
\draw[gray] (-5,-.55)--(5,-.55);
\draw[thick] (-2.5,0)--(2.5,0);
\foreach \x/\i in {-2.5/1,2.5/2}{
\draw[body] (\x-.65,-.5) rectangle (\x+.65,.5);
\node at (\x,-.25){$m_{\i}$};
\draw[force] (\x,0)--(\x,1.55) node[above]{$\vec N_{\i}$};
\draw[force] (\x,0)--(\x,-1.35) node[below]{$m_{\i}\vec g$};
\draw[force] (\x,-.48)--(\x-1.45,-.48) node[below left]{$\vec F_{\text{тр}\i}$};}
\draw[force] (-2.5,0)--(-.95,0) node[above]{$\vec T_1$};
\draw[force] (2.5,0)--(.95,0) node[above]{$\vec T_2$};
\draw[force] (2.5,0)--(4.55,0) node[above]{$\vec F$};
\draw[motion] (-.4,1.5)--(1.1,1.5) node[right]{$\vec v$};
\end{tikzpicture}
\end{center}
Поскольку $N_1=m_1g$ и $N_2=m_2g$, запишем проекции на ось $Ox$ (вправо):
\[\begin{cases} T-\mu_1m_1g=m_1a_x,\\ F-T-\mu_2m_2g=m_2a_x.\end{cases}\]
Сложим уравнения: внутренние силы натяжения сокращаются.
\[a_x=\frac{F-\mu_1m_1g-\mu_2m_2g}{m_1+m_2},\qquad T=m_1a_x+\mu_1m_1g.\]
\example{Пример. Ускорение и натяжение нити.}
$m_1=2\unit{кг}$, $m_2=3\unit{кг}$, $F=25\unit{Н}$, $\mu_1=\mu_2=0,2$, $g=10\unit{м/с}^2$.
\[F_{\text{тр}1}=4\unit{Н},\quad F_{\text{тр}2}=6\unit{Н},\]
\[a_x=\frac{25-4-6}{2+3}=3\unit{м/с}^2,\qquad T=2\cdot3+4=10\unit{Н}.\]
Проверка по второму бруску: $(25-10-6)\unit{Н}=9\unit{Н}=3\unit{кг}\cdot3\unit{м/с}^2$. Полученное натяжение положительно, что согласуется с условием о натянутой нити.
\clearpage
\section{Движение вдоль вертикальной стены}
Брусок скользит вверх по вертикальной стене, расположенной слева. Приложенная сила $F$ направлена вверх и к стене под углом $\alpha$ к вертикали ($0<\alpha<90^\circ$). Выберем оси как на схеме: $Ox$ вправо, $Oy$ вниз.
\begin{center}
\begin{tikzpicture}[scale=1]
\draw[gray,thick] (0,-2.05)--(0,2.05);
\foreach \y in {-2,-1.75,...,2}{\draw[gray] (0,\y)--(-.18,\y-.18);}
\draw[axis] (-2.4,0)--(3.5,0) node[right]{$x$};
\draw[axis] (.45,1.95)--(.45,-2.15) node[below]{$y$};
\draw[body] (0,-.5) rectangle (.9,.5);
\coordinate (O) at (.45,0);
\coordinate (P) at (-1.25,1.65);
\draw[force] (O)--(P) node[above left]{$\vec F$};
\draw[gray,dashed] (P)--(.45,1.65) -- (O) --(-1.25,0)--(P);
\draw[force,densely dashed] (O)--(-1.25,0) node[below left]{$F\sin\alpha$};
\draw[force,densely dashed] (O)--(.45,1.65) node[above right]{$F\cos\alpha$};
\draw[accent] (.45,.78) arc[start angle=90,end angle=135.85,radius=.78] node[pos=.55,above left]{$\alpha$};
\draw[force] (O)--(2.4,0) node[above]{$\vec N$};
\draw[force] (.38,0)--(.38,-1.22) node[below left,xshift=-3pt]{$m\vec g$};
\draw[force] (.57,0)--(.57,-1.68) node[right]{$\vec F_{\text{тр}}$};
\draw[motion] (3.05,-.75)--(3.05,1.25) node[above]{$\vec v$};
\end{tikzpicture}
\end{center}
Горизонтальная составляющая силы направлена влево и прижимает брусок к стене. При отсутствии горизонтального ускорения
\[N-F\sin\alpha=0,\qquad \fr=\mu F\sin\alpha.\]
Вертикальная составляющая тяги направлена вверх. Сила тяжести и трение направлены вниз:
\[mg+\mu F\sin\alpha-F\cos\alpha=ma_y.\]
Если обозначить проекцию ускорения на направление вверх через $a_{\uparrow}=-a_y$, то
\[\boxed{a_{\uparrow}=\frac{F\cos\alpha-mg-\mu F\sin\alpha}{m}.}\]
\example{Пример.} $m=2\unit{кг}$, $F=40\unit{Н}$, $\alpha=60^\circ$, $\mu=0,1$, $g=10\unit{м/с}^2$. Тогда $N=20\sqrt3\unit{Н}$, $\fr=2\sqrt3\unit{Н}$, $F\cos60^\circ=20\unit{Н}$ и $a_{\uparrow}=-\sqrt3\unit{м/с}^2$. Брусок движется вверх, замедляясь.
\clearpage
\section*{Алгоритм: решение задачи по динамике}
\vspace{7mm}
\begin{center}
\begin{tikzpicture}[node distance=9mm]
\node[terminal] (start) {Начало};
\node[flow,below=of start] (obj) {Выберите тело или систему тел. Запишите данные в СИ.};
\node[flow,below=of obj] (forces) {Изобразите все внешние силы. Отметьте направление скольжения.};
\node[flow,below=of forces] (ax) {Выберите оси. Определите знаки проекций сил и ускорения.};
\node[flow,below=of ax] (newton) {Запишите второй закон Ньютона для каждого тела:\\ $\sum F_x=ma_x,\quad\sum F_y=ma_y$.};
\node[flow,below=of newton] (norm) {Найдите $N$ из уравнения поперёк опоры; при скольжении используйте $F_{\text{тр}}=\mu N$.};
\node[flow,below=of norm] (solve) {Решите систему относительно искомых величин. Для связанных тел учтите одинаковое ускорение вдоль нити.};
\node[diamond,draw=accent,aspect=3,align=center,inner sep=1pt,below=of solve] (check) {Знаки, единицы и\\условия согласованы?};
\node[terminal,below=of check] (end) {Запишите ответ};
\foreach \a/\b in {start/obj,obj/forces,forces/ax,ax/newton,newton/norm,norm/solve,solve/check}{\draw[link] (\a)--(\b);}
\draw[link] (check)--node[right]{\small да}(end);
\coordinate (return) at (-6.55,0);
\draw[link] (check.west)--node[above]{\small нет}(check.west -| return) |- (forces.west);
\end{tikzpicture}
\end{center}
\clearpage
\section{Самопроверка и типичные ошибки}
\subsection{Формулы, которые следует знать}
\begin{center}
\renewcommand{\arraystretch}{1.55}
\begin{tabularx}{\linewidth}{@{}l X@{}}\toprule
Формула & Когда применяется\\\midrule
$\sum\vec F=m\vec a$ & Для тела в инерциальной системе отсчёта.\\
$F_{\text{тяж}}=mg$ & Сила тяжести вблизи поверхности Земли.\\
$\fr=\mu N$ & При скольжении относительно опоры.\\
$N=mg$ & Горизонтальная опора, нет иных вертикальных сил и вертикального ускорения.\\
$N=F\sin\alpha$ & Вертикальная стена; $\alpha$ отсчитывается от вертикали, иных горизонтальных сил нет.\\
$T=m_1a_x+\mu_1m_1g$ & Для левого бруска в схеме раздела 3.\\\bottomrule
\end{tabularx}
\end{center}
\subsection{Перед записью ответа проверьте}
\begin{itemize}[label=$\square$,itemsep=7pt]
\item Все ли силы изображены именно для рассматриваемого тела?
\item Перпендикулярна ли реакция опоры поверхности?
\item Направлено ли трение против относительного скольжения?
\item Правильно ли определены знаки проекций с учётом выбранных осей?
\item От какой оси отсчитывается угол? При отсчёте от вертикали вертикальная составляющая по модулю равна $F\cos\alpha$.
\item Найдена ли реакция опоры из второго закона Ньютона перед вычислением трения?
\item При сложении уравнений двух тел исключены ли только внутренние силы натяжения?
\item Проверены ли единицы измерения, знак ускорения и условия контакта ($N\geq0$) и натяжения ($T\geq0$)?
\end{itemize}
\textcolor{darkaccent}{\textbf{Помните:}} если $a_x<0$, ускорение направлено против положительной оси $Ox$. Это само по себе не означает, что тело движется влево. При остановке модель трения скольжения требует пересмотра.
\clearpage
\section*{Тренировочный блок}
Во всех задачах примите $g=10\unit{м/с}^2$. Трение является трением скольжения; сопротивление воздуха не учитывайте. Для связанных брусков нить лёгкая, натянутая и нерастяжимая. В задачах 1--7 бруски скользят вправо по горизонтальному столу; сила тяги горизонтальна и направлена вправо.
\subsection{Базовый уровень}
\begin{enumerate}[itemsep=5pt]
\item Масса бруска $4\unit{кг}$. Найдите силу тяжести и реакцию опоры.
\item Масса бруска $6\unit{кг}$, коэффициент трения $0,25$. Найдите силу трения.
\item Масса бруска $3\unit{кг}$, тяга $21\unit{Н}$, $\mu=0,2$. Найдите ускорение и его направление.
\item Масса бруска $8\unit{кг}$, тяга $16\unit{Н}$. Брусок движется равномерно. Найдите $\mu$.
\end{enumerate}
\subsection{Средний уровень}
\begin{enumerate}[start=5,itemsep=5pt]
\item Масса бруска $5\unit{кг}$, тяга $8\unit{Н}$, $\mu=0,3$. Найдите проекцию ускорения на ось вправо. Увеличивается или уменьшается скорость?
\item Два бруска массами $2\unit{кг}$ (левый) и $4\unit{кг}$ (правый) связаны нитью. К правому приложена сила $36\unit{Н}$; для обоих $\mu=0,2$. Найдите ускорение и натяжение нити.
\item Два связанных бруска: левый массой $3\unit{кг}$, $\mu_1=0,1$; правый массой $5\unit{кг}$, $\mu_2=0,3$. К правому приложена тяга $42\unit{Н}$. Найдите ускорение и натяжение нити.
\end{enumerate}
\subsection{Уровень выше среднего}
В задачах 8--10 брусок скользит вверх вдоль стены слева. Сила $F$ направлена вверх и к стене; угол $\alpha$ отсчитывается от вертикали.
\begin{enumerate}[start=8,itemsep=5pt]
\item $F=60\unit{Н}$, $\alpha=30^\circ$, $\mu=0,2$. Найдите реакцию стены и силу трения.
\item $m=2\unit{кг}$, $F=50\unit{Н}$, $\alpha=30^\circ$, $\mu=0,4$. Найдите ускорение (проекцию на ось вверх и направление).
\item $m=3\unit{кг}$, $\alpha=45^\circ$, $\mu=0,25$. Какой должна быть сила $F$, чтобы брусок скользил вверх равномерно? Найдите также реакцию стены.
\end{enumerate}
\clearpage
\section*{Ответы}
\begin{center}
\renewcommand{\arraystretch}{1.55}
\begin{tabularx}{\linewidth}{@{}r X@{}}\toprule
№ & Ответ\\\midrule
1 & $F_{\text{тяж}}=40\unit{Н}$; $N=40\unit{Н}$.\\
2 & $F_{\text{тр}}=15\unit{Н}$.\\
3 & $a=5\unit{м/с}^2$, вправо.\\
4 & $\mu=0,2$.\\
5 & $a_x=-1,4\unit{м/с}^2$; скорость уменьшается.\\
6 & $a=4\unit{м/с}^2$, вправо; $T=12\unit{Н}$.\\
7 & $a=3\unit{м/с}^2$, вправо; $T=12\unit{Н}$.\\
8 & $N=30\unit{Н}$; $F_{\text{тр}}=6\unit{Н}$.\\
9 & $a_{\uparrow}=12,5\sqrt3-15\approx6,65\unit{м/с}^2$, вверх.\\
10 & $F=40\sqrt2\approx56,57\unit{Н}$; $N=40\unit{Н}$.\\\bottomrule
\end{tabularx}
\end{center}
\vspace{10mm}
\begin{center}{\large\color{darkaccent} Итоговая памятка}\end{center}
\renewcommand{\arraystretch}{1.3}
\begin{tabularx}{\linewidth}{@{}>{\raggedright\arraybackslash}p{4.8cm} X@{}}\toprule
Сигнал в задаче & Полезное действие\\\midrule
Тело на опоре & Сначала запишите уравнение поперёк опоры.\\
Скольжение & Определите направление трения; найдите $\mu N$.\\
Нить между телами & Запишите уравнения отдельно, затем сложите их для поиска ускорения.\\
Сила под углом & Разложите силу на две перпендикулярные составляющие.\\
Равномерное движение & Положите ускорение равным нулю.\\
Отрицательная проекция & Проверьте направление ускорения относительно оси.\\\bottomrule
\end{tabularx}
\end{document}
